\subsection{Storage}\label{sec:upstorage}

This replaces \tabref{tab:datasetSizingOps}, and the USDF totals, Qserv and
tape rows of \tabref{tab:storageFloorOps}.

\subsubsection{Bytes written and on-floor need}

A data-release campaign writes far more than it keeps. Most of its output is
intermediate data (warps, templates, pre-detection images and footprints),
created and deleted during the run. What has to be bought is the disk in use
at the campaign's high-water mark, plus everything retained from earlier
releases and from Prompt Processing. The model keeps the two quantities
apart: \emph{bytes written} is what a campaign produces; \emph{on-floor need}
is what storage purchases follow.

\smpara{Bytes written.}
The DR1 table with compression (\S1.2.2 of \emph{DRP Resource Usage
Estimates}) totals 47\,PB for the estimate's \smDRoneVisits{} visits. The model
sums it dataset type by dataset type and scales it by the same \smYrScale{}
used for compute, giving \smBytesWrittenPB{}\,PB per survey year. Release
products retained after the campaign (compressed \texttt{visit\_image} plus
the standardized catalogues) are \smRetainedPB{}\,PB per survey year; the other
\smIntermediatesPB{}\,PB is transient. For DP2 the same page (\S1.1.3)
estimates \smDPtwoOutPB{}\,PB written, of which \smDPtwoRelPB{}\,PB are release
products, for the $\sim$30k-visit DP2 visit list. DP2 fitted within the storage
already installed.

\smpara{On-floor need} in a loop year is the sum of:
\begin{itemize}
\item \emph{Live intermediates}: the share of the campaign's intermediates on
      disk at its high-water mark. This model takes \smLiveFrac{}. It is the
      \textbf{least well-founded and most influential} number in the storage
      model; it has not been measured and should be replaced from campaign
      telemetry.
\item \emph{Retained releases}, over a sliding window of \smRetWin{}
      releases.
\item \emph{Raw and calibration data}, cumulative: raw data is never deleted.
\item \emph{Prompt products}: Tier-1, kept long term and cumulative, plus a
      rolling Tier-2 window (RFC-1134).
\item An \emph{other and miscellaneous allowance} of \smOtherFrac{} of the
      subtotal, standing in for user areas, the APDB and simulation output,
      which are not itemised. It is shown as its own line so it can be
      audited.
\end{itemize}

The need is \smDPtwoOnFloorPB{}\,PB for DP2, rising to \smDRoneOnFloorPB{}\,PB
at DR1 and \smDRnineOnFloorPB{}\,PB at DR9, against about
\smInstalledPB{}\,PB installed today (\tabref{tab:smStorage}). Most of the DR1
figure depends on the live-intermediate fraction: live intermediates alone are
\smDRoneLiveInterPB{}\,PB at DR1.

\subsubsection{Fast tier}

\tabref{tab:storageFloorOps} kept a separate fast tier for the APDB and the
Qserv Object table. The 2026 model sizes the same tier from its own APDB and
Qserv figures: \smFastDRoneTB{}\,TB at DR1, rising to \smFastDRnineTB{}\,TB at
DR9 (\tabref{tab:smQserv}). The tier is \textbf{sized but not separately costed}. Its
volume is in the forecast: the APDB within the on-floor total, costed at the
disk rate, and the Object table on the Qserv nodes. What is missing is any
premium for low-latency storage, because no current price per TB is available
for it. The split of the remaining on-floor need between normal
disk and object store is not updated.

\subsubsection{File counts}

The dataset-type table also gives file counts. A DR1-scale campaign writes
several hundred million files (\tabref{tab:smStorage}), which loads filesystem
metadata as well as capacity. Two operations-era products are not yet in the
breakdown:
\begin{itemize}
\item Object Scarlet models: about 11\,TB at Early DP2 (EDP2) scale.
\item Single-band HiPS maps: about 18\,TB in some 36 million small files.
\end{itemize}
Neither yet has a DR1-scale size.

\subsubsection{Tape}

Tape is cumulative: raw and calibration data and every data release are
preserved, so the footprint only grows. It reaches \smDRnineTapePB{}\,PB by DR9,
including archive already on the floor (\tabref{tab:smTape}). Two inputs are
\emph{unverified}:
\begin{itemize}
\item \emph{Archive already on the floor}, \smTapePreTB{}\,TB. Only part of
      this could be independently confirmed.
\item \emph{Tape capacity assumed available to Rubin}, \smTapeOwnedTB{}\,TB.
      S3DF is shared by several tenants, and this figure comes from a hedged
      reading of shared accounting. Owning media reserves neither library
      slots nor drive bandwidth. Setting it to zero gives the no-headroom
      case.
\end{itemize}

\begin{landscape}
% GENERATED by generate_model.py from lsst/rubin-sizing-model @ v2026.09.2. Do not edit by hand.
% Regenerate: git clone --branch v2026.09.2 https://github.com/lsst/rubin-sizing-model ; then, inside the clone, uv run generate_model.py --emit-tex <this directory>

\tiny \begin{longtable}{|p{0.22\textwidth}|l|r|r|r|r|r|r|r|r|r|r|}
\caption{Storage with all processing at the USDF: bytes written by each campaign against the on-floor disk that has to be bought. \label{tab:smStorage}}\\
\hline
\textbf{Metric} & \textbf{Unit} & \textbf{DP2 (actual)} & \textbf{DR1} & \textbf{DR2} & \textbf{DR3} & \textbf{DR4} & \textbf{DR5} & \textbf{DR6} & \textbf{DR7} & \textbf{DR8} & \textbf{DR9} \\
 &  & FY2026 & FY2027 & FY2028 & FY2029 & FY2030 & FY2031 & FY2032 & FY2033 & FY2034 & FY2035 \\
\hline\endfirsthead
\hline
\textbf{Metric} & \textbf{Unit} & \textbf{DP2 (actual)} & \textbf{DR1} & \textbf{DR2} & \textbf{DR3} & \textbf{DR4} & \textbf{DR5} & \textbf{DR6} & \textbf{DR7} & \textbf{DR8} & \textbf{DR9} \\
 &  & FY2026 & FY2027 & FY2028 & FY2029 & FY2030 & FY2031 & FY2032 & FY2033 & FY2034 & FY2035 \\
\hline\endhead
Campaign output this FY — BYTES WRITTEN & PB & 7.4 & 52.5 & 105.0 & 157.5 & 210.0 & 262.5 & 315.0 & 367.5 & 420.0 & 472.5 \\ \hline
\quad of which release products (retained) & PB & 0.3 & 1.8 & 3.6 & 5.4 & 7.1 & 8.9 & 10.7 & 12.5 & 14.3 & 16.1 \\ \hline
\quad of which intermediates (transient) & PB & 7.1 & 50.7 & 101.4 & 152.2 & 202.9 & 253.6 & 304.3 & 355.0 & 405.8 & 456.5 \\ \hline
On floor: live intermediates (peak) & PB & 2.5 & 17.8 & 35.5 & 53.3 & 71.0 & 88.8 & 106.5 & 124.3 & 142.0 & 159.8 \\ \hline
On floor: retained releases (sliding window) & PB & 0.1 & 1.8 & 5.4 & 10.7 & 16.1 & 21.4 & 26.8 & 32.1 & 37.5 & 42.8 \\ \hline
On floor: raw + calibration, cumulative & PB & 1.1 & 2.2 & 3.4 & 4.5 & 5.6 & 6.7 & 7.8 & 8.9 & 10.1 & 11.2 \\ \hline
On floor: prompt products (Tier-1 cumulative + Tier-2) & PB & 1.9 & 3.5 & 5.2 & 6.9 & 8.6 & 10.3 & 11.9 & 13.6 & 15.3 & 17.0 \\ \hline
Subtotal of modelled datasets & PB & 5.6 & 25.3 & 49.4 & 75.3 & 101.2 & 127.1 & 153.0 & 178.9 & 204.8 & 230.7 \\ \hline
Other / misc allowance (user areas, APDB, sims) & PB & 1.1 & 5.1 & 9.9 & 15.1 & 20.2 & 25.4 & 30.6 & 35.8 & 41.0 & 46.1 \\ \hline
ON-FLOOR DISK NEED (drives purchases) & PB & 6.7 & 30.4 & 59.3 & 90.4 & 121.5 & 152.6 & 183.6 & 214.7 & 245.8 & 276.9 \\ \hline
Installed storage today & PB & 30 & 30 & 30 & 30 & 30 & 30 & 30 & 30 & 30 & 30 \\ \hline
Files written by this campaign & millions & 123 & 863 & 1,727 & 2,590 & 3,454 & 4,317 & 5,181 & 6,044 & 6,908 & 7,771 \\ \hline
Files resident on floor & millions & 46 & 321 & 643 & 964 & 1,285 & 1,606 & 1,928 & 2,249 & 2,570 & 2,891 \\ \hline
\end{longtable} \normalsize

% GENERATED by generate_model.py from lsst/rubin-sizing-model @ v2026.09.2. Do not edit by hand.
% Regenerate: git clone --branch v2026.09.2 https://github.com/lsst/rubin-sizing-model ; then, inside the clone, uv run generate_model.py --emit-tex <this directory>

\tiny \begin{longtable}{|p{0.22\textwidth}|l|r|r|r|r|r|r|r|r|r|r|}
\caption{Cumulative tape footprint and purchases. Archived data is never deleted. \label{tab:smTape}}\\
\hline
\textbf{Metric} & \textbf{Unit} & \textbf{DP2 (actual)} & \textbf{DR1} & \textbf{DR2} & \textbf{DR3} & \textbf{DR4} & \textbf{DR5} & \textbf{DR6} & \textbf{DR7} & \textbf{DR8} & \textbf{DR9} \\
 &  & FY2026 & FY2027 & FY2028 & FY2029 & FY2030 & FY2031 & FY2032 & FY2033 & FY2034 & FY2035 \\
\hline\endfirsthead
\hline
\textbf{Metric} & \textbf{Unit} & \textbf{DP2 (actual)} & \textbf{DR1} & \textbf{DR2} & \textbf{DR3} & \textbf{DR4} & \textbf{DR5} & \textbf{DR6} & \textbf{DR7} & \textbf{DR8} & \textbf{DR9} \\
 &  & FY2026 & FY2027 & FY2028 & FY2029 & FY2030 & FY2031 & FY2032 & FY2033 & FY2034 & FY2035 \\
\hline\endhead
Raw + calibration archived, cumulative & PB & 1.1 & 2.2 & 3.4 & 4.5 & 5.6 & 6.7 & 7.8 & 8.9 & 10.1 & 11.2 \\ \hline
Release products archived, cumulative & PB & 0.3 & 1.8 & 5.4 & 10.7 & 17.8 & 26.8 & 37.5 & 50.0 & 64.2 & 80.3 \\ \hline
Backup of on-floor data & PB & 0.0 & 0.0 & 0.0 & 0.0 & 0.0 & 0.0 & 0.0 & 0.0 & 0.0 & 0.0 \\ \hline
TAPE FOOTPRINT, cumulative (incl. archive already on floor) & PB & 9.6 & 12.2 & 16.9 & 23.4 & 31.6 & 41.7 & 53.5 & 67.1 & 82.5 & 99.7 \\ \hline
Tape capacity assumed available (UNVERIFIED) & TB & 66,300 & 66,300 & 66,300 & 66,300 & 66,300 & 66,300 & 66,300 & 66,300 & 66,300 & 66,300 \\ \hline
Cumulative tape capacity to buy & TB & 0 & 0 & 0 & 0 & 0 & 0 & 0 & 788 & 16,177 & 33,350 \\ \hline
Tape purchased this FY & TB & 0 & 0 & 0 & 0 & 0 & 0 & 0 & 788 & 15,389 & 17,173 \\ \hline
\end{longtable} \normalsize

\end{landscape}
